\section{Introduction}

Business rules are operational rules that organizations follow to carry out their activities\cite{Huan96a}. They define the core logic of business processes and ensure that these processes remain consistent with organizational policies, objectives, and regulations. When implemented in software systems, business rules play a critical role in supporting daily operations\cite{Cose13a}. They can range from simple validation rules to complex decision-making mechanisms.

However, business rules are rarely implemented as a single, clearly identifiable component of an Information System (IS). Instead, they are often scattered across different parts of the source code\cite{Cose12a}. In long-lived software systems, these rules are also affected by the evolution of the codebase. Systems are developed and maintained over many years by different developers and teams, making their implementations increasingly difficult to understand. Documentation may become incomplete or outdated, while the developers who originally implemented the rules may no longer be available\cite{Cose13a}. As a result, understanding how a particular business rule is implemented can become a difficult task.

At the same time, the rules governing organizations continuously evolve. Changes in these rules may result from new market conditions, organizational needs, or regulatory requirements\cite{Chit19a}. Such changes often require corresponding modifications to the software. Before making a change, developers must identify the parts of the source code that implement the affected business rule. In a large and complex system, this task can be difficult and time-consuming. Developers may spend considerable effort searching for the relevant code before they can even begin implementing the change. This process is also error-prone, as modifying the wrong parts of the code can lead to inconsistent behavior or unintended changes to the system.

This challenge also arises in industrial software systems, where business rules are implemented in large and evolving codebases. In this paper, we investigate this problem through a real-world case provided by Berger-Levrault, a company that develops information systems for public administration, healthcare, and education.

In this paper, we propose an approach to assist developers in locating the source code that implements a given business rule. Given a textual description of a business rule, our approach identifies the methods that are most likely to implement it and progressively reduces the set of candidate methods. This reduces the amount of code that developers need to inspect and allows them to focus their investigation on the parts of the system that are most relevant to the rule.

A further challenge concerns the evaluation of such an approach. In real-world software systems, there is often no ground truth, that is, no authoritative mapping between business rules and the code that implements them. Without such a reference, directly measuring the effectiveness of a business rule localization approach is difficult. In our industrial case, a small number of business rules have explicit links to their implementations. We use these links as ground truth for evaluating our approach.

The overall research question addressed in this work is:

\textbf{How can we accurately identify the implementation of a specific business rule from its textual description?}

To investigate this question, we decompose it into two research questions focusing on candidate identification and candidate reduction:

\begin{enumerate}
    \item[\textbf{RQ1.}] To what extent can our approach identify candidate methods for a given business rule while retaining its correct implementation?
    \item[\textbf{RQ2.}] To what extent can information from business-rule documentation further reduce the set of candidate methods after applying our approach?
\end{enumerate}

The remainder of this paper is organized as follows. Section~\ref{sec:background} introduces the main concepts underlying business rules and their representation in source code, while Section~\ref{sec:relatedWork} reviews related work on business-rule identification and localization. Section~\ref{sec:industrialContext} presents the industrial context and problem setting that motivate our study. Section~\ref{sec:approach} describes the proposed approach for addressing RQ1, including its main steps for identifying candidate methods. Section~\ref{sec:evaluation} presents the evaluation strategy and the rule--implementation links used as ground truth, followed by the case study in Section~\ref{sec:caseStudy}. Section~\ref{sec:resultsRQ1} reports and analyzes the results obtained for RQ1. Section~\ref{sec:useDocumentation} then introduces the use of business-rule documentation as an additional source of contextual information to address RQ2, and Section~\ref{sec:resultsRQ2} evaluates its contribution to further refining the candidate sets. Section~\ref{sec:discussion} discusses the main findings and their implications, while Section~\ref{sec:threats} discusses the threats to validity. Section~\ref{sec:futureWork} outlines directions for future work. Finally, Section~\ref{sec:conclusion} concludes the paper.